% Version fixed128 standalone insertion, fully integrated into the mother TeX.
% New non-pro-2 residual S_3 sector, integral Schreier transfer, Hecke
% projector, and a coefficient-specific full marked Jacobian Smith certificate.

\subsection{The first non-pro-$2$ residual sector: tame $S_3$, Schreier,
Hecke, and a native integral Jacobian certificate}
\label{subsec:dyadic-s3-schreier-hecke}

The strict Nielsen bridge of
Theorem~\ref{thm:dyadic-tame-jacobian-fox-splitting}
assumes that the \emph{entire} residual image is pro-$2$.
The smallest honest test outside that domain is the tame
$S_3=\mathrm{GL}_2(\mathbf F_2)$ quotient of $G_{\mathbf Q_2}$.
For this quotient the Sylow preimage has odd index three
but is nonnormal.  We give its actual coset words,
integral transfer and Hecke identities, and a nontrivial
second-degree coefficient calculation.  The comparison is
at the coefficient level; no uniform full-group
resolution is inferred from the two marked relators.

Put $\Gamma=G_{\mathbf Q_2}$, choose
$\alpha^3=2$, and define
\[
F=\mathbf Q_2(\alpha),\qquad
L=F(\zeta_3).
\]
Let $\pi:\Gamma\twoheadrightarrow S_3$ be the
permutation action on the three cubic roots,
and write $N=\ker\pi$.  Choose compatible tame
Roe--Turturean generators $\sigma,\tau$ whose
images are a transposition $s$ and a $3$-cycle $u$:
$sus^{-1}=u^{-1}=u^2$.
The wild marked generators $x_0,x_1$
map to $1$ in this tame quotient.
Take the Sylow subgroup $S=\langle s\rangle$
and $H=\pi^{-1}(S)=G_F$, so that
\[
[\Gamma:H]=3,\qquad [H:N]=2.
\]

\begin{theorem}[A concrete non-pro-$2$ integral residual carrier]
\label{thm:dyadic-s3-integral-sector}
The extension $F/\mathbf Q_2$ is totally
and tamely ramified of degree $3$;
$L/F$ is unramified quadratic, and
$\operatorname{Gal}(L/\mathbf Q_2)\simeq S_3$.
On the standard integral rank-two
representation $T=\mathbf Z_2^2$, the
chosen images act by
\begin{equation}\label{eq:dyadic-s3-integral-matrices}
\boxed{\quad
\rho_T(\sigma)=
S_0=\begin{pmatrix}0&1\\1&0\end{pmatrix},
\qquad
\rho_T(\tau)=
U_0=\begin{pmatrix}-1&-1\\1&0\end{pmatrix},
\qquad
\rho_T(x_0)=\rho_T(x_1)=I_2.
\quad}
\end{equation}
They satisfy $S_0^2=U_0^3=I_2$ and
$S_0U_0S_0^{-1}=U_0^2$.
Reduction modulo $2$ gives the faithful
natural representation
$\bar\rho_T:\Gamma\twoheadrightarrow
\mathrm{GL}_2(\mathbf F_2)$,
whose image is \emph{not} a $2$-group,
while $\rho_T(H)=\langle S_0\rangle$
is a $2$-group.
The maximal pro-$2$ local quotient
$P_H=G_F(2)$ is a Demu\v{s}kin group
of rank five with $q(P_H)=2$ and
full cyclotomic orientation image
$\mathbf Z_2^\times$.
It therefore admits a (not yet
Schreier-marked) five-generator
one-relator model
\begin{equation}\label{eq:dyadic-s3-h-demushkin-five-relator}
P_H\cong\bigl\langle z_1,\ldots,z_5
\,\big|\,
z_1^2z_2^4[z_2,z_3][z_4,z_5]
\bigr\rangle_{\mathrm{pro}-2}.
\end{equation}

The complete integral cohomology ledgers
in degrees zero through two are
\begin{equation}\label{eq:dyadic-s3-natural-cohom-ledger}
\boxed{
\begin{array}{c|ccc}
 &H^0(-,T)&H^1(-,T)&H^2(-,T)\\ \hline
 \Gamma&0&\mathbf Z_2^{\,2}&0\\
 H&\mathbf Z_2&\mathbf Z_2^{\,7}&\mathbf F_2 .
\end{array}}
\end{equation}
Consequently the integral Sylow transfer
idempotent
\begin{equation}\label{eq:dyadic-s3-sylow-idempotent}
e_H=\tfrac13\,\operatorname{res}_H^\Gamma
\operatorname{cor}_H^\Gamma
\end{equation}
annihilates $H^0(H,T)$ and $H^2(H,T)$.
On $H^1(H,T)\cong\mathbf Z_2^7$
it is an actual integral projection
onto a free rank-two summand, with
free rank-five kernel.
\end{theorem}

\begin{proof}
The polynomial $X^3-2$ is Eisenstein
at $2$, so $F/\mathbf Q_2$ has degree
three and ramification index three.
Because this index is odd, the
ramification is tame.  The polynomial
$X^2+X+1$ is irreducible modulo $2$;
its root $\zeta_3$ generates the
unramified quadratic extension.
The Galois closure is thus a
six-element tame extension with the
usual $S_3$ action.
The matrices in
\eqref{eq:dyadic-s3-integral-matrices}
satisfy the displayed finite-group
relations, and their reductions
generate $\mathrm{GL}_2(\mathbf F_2)$.
They define the standard integral
permutation-difference lattice.
The odd-degree extension $F/\mathbf Q_2$
contains no $\mu_4$, so the exact
$2$-primary roots of unity in $F$
form $\mu_2$.  Moreover the
cyclotomic pro-$2$ tower of
$\mathbf Q_2$ has trivial
intersection with $F$, since
their finite intersection would
have degree dividing both $3$
and a power of $2$.
The orientation image for $G_F(2)$
is therefore all of $\mathbf Z_2^\times$.
The local Demu\v{s}kin classification
then gives rank $[F:\mathbf Q_2]+2=5$,
$q=2$, and the odd-rank $f=2$
standard relation
\eqref{eq:dyadic-s3-h-demushkin-five-relator}.

For the cohomology computation put
$V=T\otimes_{\mathbf Z_2}\mathbf Q_2$ and
$D=T^\vee(1)\otimes_{\mathbf Z_2}
(\mathbf Q_2/\mathbf Z_2)$.
The wild inertia $W$ acts trivially
on $T$ because $\rho_T$ is tame.
Its cyclotomic image contains $-1$
(and in fact the entire cyclotomic
pro-$2$ image), so
$D^W=D[2]=T^\vee/2T^\vee$.
The finite residual $S_3$ action
on $T^\vee/2$ has no nonzero
invariants: the $3$-cycle $U_0$
has no eigenvalue $1$ modulo $2$.
Integral local Tate duality
therefore gives $H^2(\Gamma,T)=0$.
For $H$, the representation is
unramified quadratic, and
$(T^\vee/2)^{\langle S_0\rangle}$
is one-dimensional, giving
$H^2(H,T)\cong\mathbf F_2$.

The fixed lattices are
$T^\Gamma=0$ and
$T^H=\mathbf Z_2(1,1)^t$.
The local Euler--Poincar\'e rank
formula \cite{NSW} gives
\[
\operatorname{rank}_{\mathbf Z_2}H^1(\Gamma,T)=2,
\qquad
\operatorname{rank}_{\mathbf Z_2}H^1(H,T)=2[ F:\mathbf Q_2]+1=7.
\]
These groups have no $2$-primary
torsion.  Indeed the invariants in
$V/T$ are zero for $\Gamma$, since
$(T/2T)^\Gamma=0$; for $H$ the
invariants $(V/T)^H$ are the full
diagonal $\mathbf Q_2/\mathbf Z_2$
and are already the image of $V^H$.
The long exact sequence for
$0\to T\to V\to V/T\to0$
therefore introduces no torsion in
$H^1$.
This proves
\eqref{eq:dyadic-s3-natural-cohom-ledger}.
Finally $3$ is a unit in
$\mathbf Z_2$, and the standard
restriction--corestriction identity
makes \eqref{eq:dyadic-s3-sylow-idempotent}
a derived idempotent whose image is
$R\Gamma(\Gamma,T)$.
Its actions in degrees zero, one
and two follow from the ledger
and the splitting of finitely
generated free $\mathbf Z_2$-modules.
\end{proof}

\begin{theorem}[Finite nonnormal Schreier table and corestriction]
\label{thm:dyadic-s3-schreier-corestriction}
Choose right-coset representatives
$t_0=1,t_1=\tau,t_2=\tau^2$
for $H\backslash\Gamma$.
Write
\[
t_i g=h_{i,g}\,t_{j(i,g)},\qquad
h_{i,g}\in H,
\]
where $j(i,g)$ is determined by the
right $S_3$-coset action.
Then the complete Schreier data
for the four marked generators are
\begin{equation}\label{eq:dyadic-s3-schreier-table}
\boxed{
\begin{array}{c|ccc}
g& i=0&i=1&i=2\\ \hline
\sigma&
(0,\sigma)&(2,\sigma)&(1,\sigma\tau^3)\\
\tau&
(1,1)&(2,1)&(0,\tau^3)\\
x_j&
(0,x_j)&(1,\tau x_j\tau^{-1})&
(2,\tau^2x_j\tau^{-2}),\quad j=0,1 .
\end{array}}
\end{equation}
An entry $(k,h)$ means $j(i,g)=k$
and $h_{i,g}=h$.
In particular a finite profinite
Schreier generating system for $H$
consists of
\[
\sigma,\ \tau^3,\quad
\tau^i x_j\tau^{-i}
\quad(i=0,1,2,\ j=0,1).
\]
There are eight displayed
generators after removal of trivial
and repeated Schreier words.
The six wild conjugates have
pro-$2$ normal closure, while the
tame element $\tau^3$ has pro-odd
order.  Thus $\tau^3$ maps to $1$
under $H\twoheadrightarrow P_H=H(2)$.

For any continuous $H$-cocycle
$c\in Z^1(H,T)$, the corestriction
is represented on marked generators
by the finite coset formula
\begin{equation}\label{eq:dyadic-s3-transfer-one}
(\operatorname{cor}_H^\Gamma c)(g)
=\sum_{i=0}^{2}
  \rho_T(t_i)^{-1}c(h_{i,g}).
\end{equation}
For the specific rank-two $S_3$
module \eqref{eq:dyadic-s3-integral-matrices}
this simplifies \emph{exactly} to
\begin{equation}\label{eq:dyadic-s3-transfer-simplified}
\boxed{
\begin{aligned}
(\operatorname{cor}c)(\sigma)&=0,\\
(\operatorname{cor}c)(\tau)&=0,\\
(\operatorname{cor}c)(x_j)
 &=\sum_{i=0}^{2} U_0^{-i}
   c(\tau^i x_j\tau^{-i}),
 \qquad j=0,1 .
\end{aligned}}
\end{equation}
This is an integral formula, with
no average or division by $2$;
the normalized Sylow projector
divides by the odd index $3$ only
\emph{after} corestriction.

The analogous higher-degree bar
corestriction is obtained by tracking
successive coset indices:
\begin{equation}\label{eq:dyadic-s3-transfer-all-arity}
(\operatorname{cor}c)(g_1,\dots,g_n)
=\sum_{i=0}^2\rho_T(t_i)^{-1}
c\!\left(h_{i,g_1},
h_{j(i,g_1),g_2},\ldots,
h_{j(i,g_1\cdots g_{n-1}),g_n}\right).
\end{equation}
It is compatible with the normalized
bar differential and hence realizes
the usual corestriction on
cohomology; this is \emph{not}
by itself a finite Demu\v{s}kin
Fox change-of-basis formula.
\end{theorem}

\begin{proof}
The tame relation is
$\sigma^{-1}\tau\sigma=\tau^2$,
so $\tau^i\sigma=\sigma\tau^{2i}$
as an exact profinite word identity.
Reducing $2i$ modulo $3$ against
the chosen coset representatives
gives the first line of
\eqref{eq:dyadic-s3-schreier-table};
the second and third lines follow
from right multiplication by
$\tau$ and by a wild generator
with trivial residual $S_3$ image.
The Schreier rewriting theorem for
finite-index open profinite subgroups
gives the displayed generating
system; the normal closure of the
wild conjugates is contained in
the marked wild pro-$2$ subgroup,
and it is pro-$2$.
The element $\tau^3$ lies in tame
inertia, of pro-odd order, so its
image in the maximal pro-$2$
quotient of $H$ is trivial.

The standard bar corestriction for
right-coset representatives is
\eqref{eq:dyadic-s3-transfer-all-arity}.
For degree one the crossed-cocycle
identity and
$h_{i,g_1g_2}=
h_{i,g_1}h_{j(i,g_1),g_2}$
show directly that its output is
a $\Gamma$-cocycle, proving
\eqref{eq:dyadic-s3-transfer-one}.
The same identity gives the bar
chain-map formula in every degree.

Because the action of $\tau^3$
on $T$ is trivial and its closed
cyclic group has pro-odd order,
every continuous crossed
homomorphism $H\to T$ vanishes on
$\tau^3$.  Therefore the
$\sigma$-column of the table gives
\[
(\operatorname{cor}c)(\sigma)
=(I+U_0^{-1}+U_0^{-2})c(\sigma)=0,
\]
using $I+U_0+U_0^2=0$.
The $\tau$ column gives zero, and
the two wild columns give exactly
\eqref{eq:dyadic-s3-transfer-simplified}.
\end{proof}

\begin{theorem}[The index-three integral Hecke--Sylow polynomial]
\label{thm:dyadic-s3-hecke-sylow-polynomial}
On $R\Gamma(H,T)$ define
\begin{equation}\label{eq:dyadic-s3-hecke-operator}
\mathsf A=\operatorname{res}_H^\Gamma
\operatorname{cor}_H^\Gamma,
\qquad
\mathsf T_{\rm Hk}=\mathsf A-I.
\end{equation}
These operators obey, as derived
endomorphisms (and hence in every
cohomological degree),
\begin{equation}\label{eq:dyadic-s3-hecke-minpoly}
\boxed{\qquad
\mathsf A^2=3\mathsf A,\qquad
\mathsf T_{\rm Hk}^2=
\mathsf T_{\rm Hk}+2I,\qquad
e_H=\frac{I+\mathsf T_{\rm Hk}}3.
\qquad}
\end{equation}
The nonidentity term also has
the exact Mackey double-coset
description
\begin{equation}\label{eq:dyadic-s3-double-coset-hecke}
\mathsf T_{\rm Hk}=
\operatorname{cor}_N^H\circ
(\tau)_*\circ
\operatorname{res}_N^H,
\qquad N=\ker\pi,
\end{equation}
where $(\tau)_*$ is conjugation on
the normal subgroup $N$, together
with its prescribed action on $T$.
Both the restriction and the
corestriction through $N$ are
integrally defined: no division
by $[H:N]=2$ occurs.

For the coefficient lattice $T$,
the two eigensummands selected by
$e_H$ have the following precise
meaning:
\[
\begin{array}{c|ccc}
& H^0(H,T)&H^1(H,T)&H^2(H,T)\\ \hline
\operatorname{im}e_H&0&\mathbf Z_2^2&0\\
\ker e_H&\mathbf Z_2&\mathbf Z_2^5&\mathbf F_2.
\end{array}
\]
Thus the genuine non-pro-$2$
descent kills a torsion top Fox
class even though the index-three
operator itself remains integral.
\end{theorem}

\begin{proof}
The standard cohomological identity
$\operatorname{cor}\operatorname{res}
=3I$ implies
$\mathsf A^2=3\mathsf A$.
Substituting $\mathsf T_{\rm Hk}=\mathsf A-I$
gives the displayed quadratic
polynomial and projector formula.
In the finite residual quotient,
$S=\langle s\rangle$ has precisely
two double cosets in $S_3$:
$S$ and $SuS$.  The latter
has intersection
$S\cap uSu^{-1}=\{1\}$,
whose inverse image in $\Gamma$
is exactly $N$.
The Mackey restriction--corestriction
formula therefore yields
\eqref{eq:dyadic-s3-double-coset-hecke}.
All maps are defined before
inverting $2$ and require no
normality of $H$.
The final table follows from
Theorem~\ref{thm:dyadic-s3-integral-sector}
and the idempotent's identification
with $R\Gamma(\Gamma,T)$.
\end{proof}

\begin{theorem}[The first non-pro-$2$ integral adjoint cohomology ledger]
\label{thm:dyadic-s3-adjoint-local-cohomology}
Retain the integral $S_3$ lattice $T$
and put
\[
M=\operatorname{ad}T
=\operatorname{End}_{\mathbf Z_2}(T),
\qquad \bar M=M/2M.
\]
The cohomology of $\Gamma=G_{\mathbf Q_2}$
is
\begin{equation}\label{eq:dyadic-s3-adjoint-integral-ledger}
\boxed{\quad
H^0(\Gamma,M)=\mathbf Z_2,\qquad
H^1(\Gamma,M)=\mathbf Z_2^{\,5},\qquad
H^2(\Gamma,M)=\mathbf F_2.
\quad}
\end{equation}
Its mod-$2$ cohomology dimensions are
\begin{equation}\label{eq:dyadic-s3-adjoint-mod2-ledger}
\boxed{\quad
\bigl(\dim H^0,\dim H^1,\dim H^2\bigr)
(\Gamma,\bar M)=(1,6,1).
\quad}
\end{equation}
The tame order-three averaging operator
is integral on both $T$ and $M$:
\begin{equation}\label{eq:dyadic-s3-tame-averaging}
\Pi_{\rm tame}^T
=\tfrac13(I+U_0+U_0^2)=0,\qquad
\Pi_{\rm tame}^M(A)
=\tfrac13\sum_{j=0}^2U_0^jAU_0^{-j}.
\end{equation}
The second projector has image
$\mathbf Z_2[U_0]$ of rank two;
conjugation by $S_0$ fixes exactly
the scalar rank-one submodule.
\end{theorem}

\begin{proof}
Over $\mathbf F_2$, the natural
representation of $S_3=\mathrm{GL}_2(\mathbf F_2)$
is absolutely irreducible and its
endomorphism centralizer consists
only of scalars.  Consequently
$M^\Gamma=\mathbf Z_2 I$ and
$\bar M^\Gamma=\mathbf F_2 I$.
As in Theorem~\ref{thm:dyadic-s3-integral-sector},
wild inertia acts trivially on $M$
but has cyclotomic image containing
$-1$.  Therefore the invariants in
\[
M^\vee(1)\otimes_{\mathbf Z_2}
(\mathbf Q_2/\mathbf Z_2)
\]
are precisely the $S_3$-invariants
of $M^\vee/2M^\vee$.
The trace pairing identifies the
latter with $\bar M$;
it has one-dimensional invariant
space.  Integral local Tate duality
gives $H^2(\Gamma,M)\cong\mathbf F_2$.
Rational local Euler--Poincar\'e
gives
$\operatorname{rank}H^1(\Gamma,M)
=4+\operatorname{rank}H^0(\Gamma,M)=5$.

There is no torsion in $H^1$:
if an invariant class in
$(M\otimes\mathbf Q_2)/M$
has denominator $2^n$, its
numerator modulo $2$ is
$S_3$-invariant and hence scalar.
Subtracting an appropriate
rational scalar matrix reduces
the denominator; iteration shows
that all these invariants come
from $(M\otimes\mathbf Q_2)^\Gamma$.
The coefficient long exact
sequence then forces
$H^1(\Gamma,M)$ to be torsion-free.
The mod-$2$ Euler formula with
$\dim H^0=\dim H^2=1$ gives
$\dim H^1=4+1+1=6$.

Finally $U_0$ has
$U_0^2+U_0+I=0$, and
$3\in\mathbf Z_2^\times$,
so the displayed tame averages
are genuine integral projectors.
The centralizer of $U_0$ in
$M_2(\mathbf Z_2)$ is
$\mathbf Z_2[U_0]$ because its
minimal polynomial is irreducible
modulo $2$.  Conjugation by
$S_0$ sends $U_0$ to $U_0^2$;
its fixed part is exactly
$\mathbf Z_2 I$.
\end{proof}

\begin{theorem}[Explicit full-marked $8\times16$ integral Jacobian certificate]
\label{thm:dyadic-s3-marked-adjoint-jacobian}
Use the full Roe--Turturean
$\Gamma$-presentation of
Theorem~\ref{thm:dyadic-q2-marked-word-jet-solver}
and the integral residual lift
\eqref{eq:dyadic-s3-integral-matrices}.
Let
\begin{equation}\label{eq:dyadic-s3-marked-adjoint-complex}
\mathcal J^\bullet_{S_3}(M):
\quad
M\xrightarrow{d^0} M^{4}
\xrightarrow{d^1} M^{2}
\end{equation}
be the completed marked two-relation
Jacobian complex at this lift.
Its differentials are the principal
cocycle differential and the first
derivatives of the literal tame/wild
marked matrix relations.
In the row-major basis
$(E_{11},E_{12},E_{21},E_{22})$
of $M$, order source variables by
$(\sigma,\tau,x_0,x_1)$
and target relation matrices by
(tame, wild).  Then the reduction
of $d^1$ modulo $2$ is the
\emph{actual} following $8\times16$
matrix:
\begin{equation}\label{eq:dyadic-s3-jacobian-mod2-matrix}
\boxed{\ \!
{\renewcommand{\arraystretch}{0.95}
\setlength{\arraycolsep}{2.5pt}
\begin{array}{rrrr|rrrr|rrrr|rrrr}
1&1&0&1&1&1&0&0&0&0&0&0&0&0&0&0\\
0&1&1&0&1&0&0&1&0&0&0&0&0&0&0&0\\
1&0&1&1&0&0&1&0&0&0&0&0&0&0&0&0\\
1&1&0&1&0&1&0&1&0&0&0&0&0&0&0&0\\ \hline
0&0&0&0&1&1&1&0&0&0&0&0&1&1&1&1\\
0&0&0&0&1&0&0&1&0&0&0&0&1&0&1&1\\
0&0&0&0&1&0&0&1&0&0&0&0&1&1&0&1\\
0&0&0&0&0&1&1&1&0&0&0&0&1&1&1&1
\end{array}}\ \! } .
\end{equation}
Its rank is $7$, and its
left nullspace is generated by
\begin{equation}\label{eq:dyadic-s3-jacobian-left-null}
\ell=(1,0,0,1\mid1,0,0,1),
\end{equation}
the sum of the two matrix traces.

More strongly, the genuine
\emph{integral} Jacobian has
Smith invariant factors
\begin{equation}\label{eq:dyadic-s3-jacobian-smith}
\boxed{\quad
\operatorname{SNF}_{\mathbf Z_2}(d^1)
=\operatorname{diag}
(1,1,1,1,1,1,1,2)
\quad}
\end{equation}
with eight additional zero
source columns.
An explicit finite-precision
certificate is obtained as follows:
work in
$\bigl(\mathbf Z/4[\varepsilon]/
(\varepsilon^2)\bigr)$, use the
left variations
\[
\rho_\varepsilon(g)
=(I+\varepsilon A_g)\rho_T(g),
\quad g\in\{\sigma,\tau,x_0,x_1\},
\]
set $\sigma^{\omega_2}=\sigma$,
and evaluate
$(x_i\tau)^{\omega_2}=(x_i\tau)^9$.
In the resulting $8\times16$
Jacobian modulo $4$, the
$8\times8$ minor formed by the
one-based columns
\begin{equation}\label{eq:dyadic-s3-jacobian-minor-columns}
\boxed{\qquad
\{1,2,5,6,9,13,14,15\}
\qquad}
\end{equation}
gives the following fully displayed
minor, with those columns in
increasing order:
\begin{equation}\label{eq:dyadic-s3-jacobian-mod4-minor}
\mathsf D_4=
\begin{pmatrix}
1&3&3&3&0&0&0&0\\
0&3&1&2&0&0&0&0\\
3&0&0&2&0&0&0&0\\
3&1&0&1&0&0&0&0\\
0&0&1&1&2&1&1&3\\
0&0&3&2&0&3&2&3\\
0&0&1&2&0&1&3&2\\
0&0&2&3&0&3&3&1
\end{pmatrix}
\pmod4,\qquad
\boxed{\det(\mathsf D_4)\equiv2\pmod4}.
\end{equation}
Together with
the rank-seven reduction
\eqref{eq:dyadic-s3-jacobian-mod2-matrix},
this proves
\eqref{eq:dyadic-s3-jacobian-smith}.

The preceding degree-zero map
has mod-$2$ rank three.
Consequently the entire marked
complex has the integral cohomology
\begin{equation}\label{eq:dyadic-s3-marked-cochain-ledger}
H^0(\mathcal J_{S_3})=\mathbf Z_2,\qquad
H^1(\mathcal J_{S_3})=\mathbf Z_2^5,\qquad
H^2(\mathcal J_{S_3})=\mathbf F_2.
\end{equation}
It is homotopy equivalent, after
finite integral row/column reduction,
to the finite elementary complex
\begin{equation}\label{eq:dyadic-s3-marked-smith-model}
\boxed{\quad
\mathbf Z_2[0]\ \oplus\
\mathbf Z_2^{\,5}[-1]\ \oplus\
\bigl[\mathbf Z_2
\xrightarrow{\ 2\ }\mathbf Z_2\bigr]_{[1,2]} .
\quad}
\end{equation}
In particular, the full marked
two-relator coefficient Jacobian
is \emph{cohomologically complete}
for this specified integral
$S_3$ \emph{adjoint} representation.
It is not asserted to be a
projective group resolution over
$\mathbf Z_2[[\Gamma]]$.
\end{theorem}

\begin{proof}
The residual group action on
$\operatorname{ad}\bar T$ is
conjugation by the two matrices
$S_0,U_0$.  In the dual-number
ring modulo $2$, every lift
of a matrix with residual order
three has $2$-primary exponent
at most two; its $\omega_2$
projection is therefore its third
power.  Modulo $4$ the corresponding
exponent divides four, so the
CRT exponent $9$, congruent to $1$
modulo $4$ and $0$ modulo $3$,
computes its $\omega_2$ projection.
For the residual order-two matrix
$S_0$, the whole lift is pro-$2$,
so its $\omega_2$ projection is
itself.  These are finite instances
of the integral binomial theorem
\ref{thm:dyadic-omega2-power-formula}.

Substitute these powers into the
literal auxiliary words
$\sigma_2,u_i,d_0,z_0,c_0,g_0,
d_g,h_c,h_0$
of Roe--Turturean, and differentiate
the two relation matrices with
respect to the sixteen standard
left tangent directions
$A_g=E_{ij}$.
This gives exactly
\eqref{eq:dyadic-s3-jacobian-mod2-matrix};
a direct row elimination yields
rank seven.
The displayed $\ell$ annihilates
every column, so the left
nullspace is exactly one-dimensional.
Performing the same calculation
modulo $4$ with exponent $9$
gives the minor specified by
\eqref{eq:dyadic-s3-jacobian-minor-columns},
of determinant $2$ modulo $4$.
Since all $8\times8$ minors
vanish modulo $2$ but at least one
has valuation exactly one,
the invariant factors are seven
units followed by a single factor
of valuation one.  This proves
\eqref{eq:dyadic-s3-jacobian-smith}.

The principal-cocycle map has
kernel equal to the scalar
centralizer $M^\Gamma=\mathbf Z_2$;
its mod-$2$ reduction has kernel
$\mathbf F_2$, so $d^0$ has
three unit invariant factors.
By $d^1d^0=0$ and integral
elementary divisors, the image
of $d^0$ is a saturated rank-three
submodule of $\ker d^1$.
Therefore $H^1$ is free of
rank $16-8-3=5$, while $H^2$
is the cokernel $\mathbf F_2$.
The degree-zero invariant
space is $\mathbf Z_2$.
Removing the seven unit
degree-one--two summands and
three unit degree-zero--one
summands proves the homotopy
normal form
\eqref{eq:dyadic-s3-marked-smith-model}.
\end{proof}

\begin{corollary}[Coefficient-specific comparison with the Sylow--Fox carrier]
\label{cor:dyadic-s3-adjoint-smith-sylow-comparison}
For the rank-four coefficient
$M=\operatorname{ad}T$ of
Theorem~\ref{thm:dyadic-s3-adjoint-local-cohomology},
both the full marked coefficient
Jacobian
$\mathcal J^\bullet_{S_3}(M)$ and
the all-Sylow Fox retract
$\operatorname{Im}(e_H)$ are bounded
perfect integral complexes
with cohomology
$(\mathbf Z_2,\mathbf Z_2^5,\mathbf F_2)$
in degrees $(0,1,2)$.
They are therefore isomorphic in
$D(\mathbf Z_2)$ and can be
connected by an integral
chain-homotopy equivalence
\emph{after choices of Smith
reduction and a top-degree
Tate-dual generator}.
The degree-zero and degree-one
identifications are canonically
fixed by invariants and evaluation
of actual framed deformations.
The top-degree identification
may be normalized by the
mod-$2$ trace obstruction
\eqref{eq:dyadic-s3-jacobian-left-null}
and local Tate duality.

This removes the \emph{existence}
and integer-divisor gap for
a direct finite coefficient
comparison in the first
non-pro-$2$ adjoint sector.
It does \emph{not} supply a
canonical strict marked-word
chain map to the chosen five-
generator Fox resolution of $P_H$,
nor does it establish a
presentation-independent
second-syzygy theorem for
all non-pro-$2$ residual sectors.
\end{corollary}

\begin{proof}
Theorem~\ref{thm:dyadic-s3-marked-adjoint-jacobian}
gives the explicit integral Smith
model of the marked complex.
The all-Sylow retract computes
$R\Gamma(\Gamma,M)$ by
Theorem~\ref{thm:all-sylow-transfer}.
Its cohomology is computed in
Theorem~\ref{thm:dyadic-s3-adjoint-local-cohomology}.
Over the discrete valuation
ring $\mathbf Z_2$, any bounded
perfect complex admits finite
elementary-divisor decompositions,
with its homotopy type determined
by its cohomology modules in each
degree.  Hence its finite
perfect representatives may be
reduced to the same elementary
complex
\eqref{eq:dyadic-s3-marked-smith-model},
which gives an integral chain
homotopy equivalence after the
chosen reductions.
The actual cocycle relation
follows from the marked
presentation and residual
pro-$2$ automaticity
(Theorem~\ref{thm:marked-residual-strictification}),
while local Tate duality selects
a nonzero generator of the
one-dimensional top torsion.
No explicit universal
Schreier-to-minimal-Fox basis is
constructed by this argument,
explaining the stated limitation.
\end{proof}

\begin{warning}[What the $S_3$ computation does not prove]
\label{warn:dyadic-s3-syzygy-scope}
The Schreier table gives finite
words in the full group, and
the all-arity bar transfer is
an exact integral operation.
Turning those ten raw
Schreier-word entries into a
canonical minimal dyadic
Demu\v{s}kin basis for
$P_H=G_{\mathbf Q_2(\sqrt[3]2)}(2)$,
with the necessary
integral relation syzygies and
an explicit multiplicative
chain comparison, is a
strictly stronger problem.
The $S_3$ natural and adjoint
coefficient results close their
own cohomology, Hecke and
Smith-normal-form ledgers,
not the general non-pro-$2$
full-word resolution problem.
The existence of the
all-Sylow derived atlas
remains unconditional and is
not upgraded here to a
uniform full-group projective
Fox resolution.
\end{warning}
